\section{Deterministic mass and the centered-moment ledger}\label{sec:ledger}

Write \(\Delta_{ij}=T_{ij}-t_0\).  The recursion implies
\[
A_0(r-j-1)+\theta(j-1)
\le\Delta_{jq}
\le\frac{A_0+\rho\theta}{1-\rho}(r-j-1)+\theta(j-1).
\tag{5.1}
\]
Consequently
\[
 a_-(r-2)\le\Delta_{jq}\le a_+(r-2).
\tag{5.2}
\]
For \(V_i=\{i+1,\ldots,r\}\), of size \(k=r-i\), put
\[
 S_i=\sum_{\{j,q\}\subset V_i}\Delta_{jq},\qquad
 R_i=\max_{j\in V_i}\sum_{q\ne j}\Delta_{jq}.
\]
Then
\[
S_i\ge a_-(r-2)\binom k2,\quad
R_i\le a_+(r-2)(k-1),\quad
\frac{R_i}{S_i}\le\frac{2a_+}{a_-k}.
\tag{5.3}
\]

\begin{lemma}[Exact accumulated multiplicities]\label{lem:multiplicity}
Let
\[
 K_r=\sum_{i=1}^{r-1}(r-i)(i-1)^2
 =\binom r3+2\binom r4.
\tag{5.4}
\]
Then
\[
 \sum_{i=1}^{r-1}S_i\ge2A_0\binom r4+\theta K_r,
 \qquad
 \sum_{1\le i<j\le r}e_{ij}^2=\theta^2K_r.
\tag{5.5}
\]
\end{lemma}

\begin{proof}
Sum the two terms of the lower bound in (5.1) with their exposure
multiplicities.  The future-degree term counts every ordered quadruple twice,
giving \(2\binom r4\); the residual term is exactly the sum in (5.4).
The second identity follows directly from \(e_{ij}=\theta(i-1)\).
\end{proof}

\subsection{Exact conditional factorization}

Fix a red-admissible \(\mathcal F_{i-1}\)-history for target \(r\) at
exposure step \(i\), and condition further on \(A_i^R\).  For \(j\in V_i\), let
\(X_j\) be the coordinate in (2.30), write \(b_j=b_{ij}\), put
\(m_j=m(b_j)\), and standardize its centered
part by
\[
 \zeta_j=\sqrt d\,(X_j-\mathbb EX_j)=\sqrt d\,X_j+m_j.
\tag{5.6}
\]
By (S2), the \(\zeta_j\) are conditionally independent, centered and
\(1\)-subgaussian.  Expanding the weighted quadratic form gives the exact
identity
\begin{align}
&\log\mathbb E\exp\left\{
-\sum_{\{j,q\}\subset V_i}T_{jq}X_jX_q\right\}\notag\\
&\qquad=-\frac1d\sum_{\{j,q\}\subset V_i}T_{jq}m_jm_q
+\Lambda_i,\tag{5.7}\\
\Lambda_i={}&\log\mathbb E\exp\left\{
\begin{aligned}
&\sum_{j\in V_i}(u_j^0+\delta u_j)\zeta_j\\
&-\frac{t_0}{d}\sum_{j<q}\zeta_j\zeta_q
-\frac1d\sum_{j<q}\Delta_{jq}\zeta_j\zeta_q
\end{aligned}\right\}.\tag{5.8}
\end{align}
where
\[
 u_j^0=\frac{t_0}{d}\sum_{q\ne j}m_q
 =\frac{m_0}{\sqrt d}\sum_{q\ne j}m_q,\qquad
 \delta u_j=\frac1d\sum_{q\ne j}\Delta_{jq}m_q.
\tag{5.9}
\]
All tilts in (5.9) are nonnegative.

\subsection{Three-factor H\"older and the square-MGF cost}

For \(k\ge2\), choose
\[
 Q_0=\frac d{4t_0k},\qquad Q_1=\frac d{4R_i},\qquad
 P^{-1}=1-Q_0^{-1}-Q_1^{-1}.
\tag{5.10}
\]
If \(R_i=0\), the increment factor is absent and we set \(Q_1=\infty\).
For all sufficiently large \(C\),
\[
 Q_0^{-1}+Q_1^{-1}
 \le\frac{4m_0}{D}+\frac{4a_+}{D^2}<\frac12,
\tag{5.11}
\]
so \(Q_0,Q_1>1\) and \(1<P\le2\).  H\"older's inequality applied to
the three random factors in (5.8) yields
\begin{align}
\Lambda_i\le{}&
\frac1P\sum_jK_{b_j}\bigl(P(u_j^0+\delta u_j)\bigr)\notag\\
&+\frac1{Q_0}\log\mathbb E
\exp\left\{-\frac{Q_0t_0}{d}\sum_{j<q}\zeta_j\zeta_q\right\}\notag\\
&+\frac1{Q_1}\log\mathbb E
\exp\left\{-\frac{Q_1}{d}\sum_{j<q}\Delta_{jq}\zeta_j\zeta_q\right\}.
\tag{5.12}
\end{align}
The middle line is the unchanged HMS uniform-quadratic factor and costs at
most \(4t_0k/d\).

For the last line, write
\(R_j=\sum_{q\ne j}\Delta_{jq}\).  Nonnegativity of \(\Delta\) gives
\[
 -\sum_{j<q}\Delta_{jq}\zeta_j\zeta_q
 \le\frac12\sum_jR_j\zeta_j^2.
\tag{5.13}
\]
Because \(Q_1R_j/(2d)\le1/8\), the pinned square-MGF estimate
\[
 \log\mathbb E e^{x\zeta_j^2}\le\frac{16}{3}x,
 \qquad 0\le x\le\frac18,
\tag{5.14}
\]
and conditional independence imply
\[
\begin{aligned}
\frac1{Q_1}\log\mathbb E
\exp\left\{-\frac{Q_1}{d}\sum_{j<q}\Delta_{jq}\zeta_j\zeta_q\right\}
&\le\frac1{Q_1}\sum_j\frac{16}{3}
\frac{Q_1R_j}{2d}\\
&=\frac{16}{3}\frac{S_i}{d},
\end{aligned}
\tag{5.15}
\]
where \(\sum_jR_j=2S_i\).  This derives the constant \(16/3\).

\subsection{The same-history CGF bridge}

Let \(P^0=P^0_{k,d}\), so that
\((P^0)^{-1}=1-Q_0^{-1}\).  The old optimized uniform ledger on this
same conditional history has deficit
\[
 \mathfrak d_i(b)=\sum_{j\in V_i}\left[
 (u_j^0)^2-\frac1{P^0}K_{b_j}(P^0u_j^0)\right].
\tag{5.16}
\]
Equations (2.6), the cutoff window, and
\(\sum_{q\ne j}m_q\ge(k-1)m_w\) give the deterministic lower bound
\[
 \mathfrak d_i(b)\ge\mathsf d_{k,w,d}.
\tag{5.17}
\]
Indeed, after writing
\(K_b(P^0u)/P^0=P^0u^2\psi_b(P^0u)\), this is exactly the
monotonicity comparison used in (2.14).  Thus the term whose accumulated
limit is \(G_*(C)\) is identified before any new factor is introduced.

It remains to compare the new linear factor with the second term in
(5.16).  Since \(K_b''\in[0,1]\), for \(u,\delta\ge0\) and \(1<P\le2\),
\[
 \frac{K_b(P(u+\delta))-K_b(Pu)}P\le2u\delta+\delta^2.
\tag{5.18}
\]
Moreover, differentiating \(K_b(Pu)/P\) and using (2.6) gives
\[
 0\le\frac{\partial}{\partial P}\frac{K_b(Pu)}P\le u^2.
\tag{5.19}
\]
As
\[
 P-P^0=PP^0\left(\frac1{P^0}-\frac1P\right)
 =\frac{PP^0}{Q_1}\le\frac{16R_i}{d},
\tag{5.20}
\]
we obtain
\begin{align}
&\frac1P\sum_jK_{b_j}\bigl(P(u_j^0+\delta u_j)\bigr)
-\frac1{P^0}\sum_jK_{b_j}(P^0u_j^0)\notag\\
&\hspace{18mm}\le
2\sum_ju_j^0\delta u_j+\sum_j(\delta u_j)^2
+\frac{16R_i}{d}\sum_j(u_j^0)^2.
\tag{5.21}
\end{align}
Each term is now bounded without probabilistic input.  From (5.3), (5.9),
\(m_j\le M_w\), and \(k/\sqrt d\le1/D\),
\begin{align}
2\sum_ju_j^0\delta u_j
&\le\frac{4m_0M_w^2}{D}\frac{S_i}{d},\tag{5.22}\\
\sum_j(\delta u_j)^2
&\le\frac{2M_w^2a_+}{D^2}\frac{S_i}{d},\tag{5.23}\\
\frac{16R_i}{d}\sum_j(u_j^0)^2
&\le\frac{32m_0^2M_w^2a_+}{a_-D^2}\frac{S_i}{d}.
\tag{5.24}
\end{align}
For completeness, (5.22) uses
\(\sum_j\delta u_j\le2M_wS_i/d\); (5.23) uses
\(\max_j\delta u_j\le M_wR_i/d\); and (5.24) uses
\[
 \sum_j(u_j^0)^2\le
 \frac{2m_0^2M_w^2}{a_-}\frac{S_i}{d},
\tag{5.25}
\]
which follows from the first inequality in (5.3) and \(r-2\ge k-1\).
The sum of (5.22)--(5.24) is precisely
\(\widehat L_w(C)S_i/d\).

Finally, at the central cutoff the deterministic \(\Delta\)-part of
(5.7) is \(-m_0^2S_i/d\).  Combining it with (5.15) and
(5.21)--(5.24) leaves
\[
 -\widehat B_w(C)\frac{S_i}{d}.
\tag{5.26}
\]
Equations (5.16)--(5.17) charge the old exact-CGF deficit once;
(5.26) charges only the new increment factor and the change
\(P^0\to P\); and (4.6) pays the boundary residual.  Hence the three
ledger entries are disjoint on the same conditional history.
